% ==========================================================================================
\section{Introduction}
\label{sec:intro}

The Lightning Network is a payment network built on Bitcoin. Its nodes run open-source software and hold their
operators' funds in payment channels, so a flaw in that software can put those funds at risk. Between July and
October~2026, every major Lightning implementation disclosed security vulnerabilities. Lightning
Labs published ten LND advisories on a single day~\cite{lightninglabs2026aug} and eight more a month
later~\cite{lightninglabs2026sep}; Core Lightning (CLN) asked operators who could not upgrade to restart their nodes
offline~\cite{coreln2026offline}; Eclair and LDK shipped security releases. Each advisory ends with the same advice:
upgrade. Whether operators do, and how fast, decides how long a disclosed vulnerability stays exploitable, and by
October CLN's maintainers reported attacks on unpatched nodes~\cite{coreln2026attacks}. For web servers, browsers and
Bitcoin and Ethereum full nodes, how fast operators patch has been measured for two decades
\cite{arbaugh2000windows,rescorla2003securityholes,frei2009firefox,durumeric2014heartbleed,kim2018ethereum}. For the
Lightning Network it has not, because Lightning gossip carries no software version: earlier studies could tell which
implementation a node runs~\cite{zabka2021nodeclassification,espinasa2024decisiontree,kumar2026clientdiversity}, not
which release.

We show that the release can be read from gossip after all, and use it to ask who patches the Lightning Network:
operators on their own schedule, the platforms that package the implementations, or the advisories themselves. Each
release announces a characteristic set of protocol feature bits. From the source of every release of LND, CLN, Eclair
and LDK since 2019 we derive which sets each release can produce, and match them against \DataUsableDays{} daily
snapshots of the public graph taken by our own node from May~2023 to October~2026. The answer is that the network
patches on release and platform schedules, not on disclosures. Most operators upgrade within months of a release;
the largest movements are platforms shipping a release to all their users at once; and an advisory, published long
after its fix, moves at most the operators who were already close to current. We make four contributions.

\begin{itemize}
\item \textbf{Release-line inference from gossip.} The release line of a public node can be read from its gossip.
A release line is the finest unit gossip separates; patch releases are invisible. Where a node states its version, the
inference is close to exact: on CLN nodes whose alias names the version, the most likely line is right for
\NVerNodeLinesExact{} of \NVerNodeLines{} node--line pairs. LND nodes state no version, so we check three things that
must hold if the inference is right, and all three do: line decreases are \PctDownLND\% of observed line changes, no
LND line released since 2024 was seen on more than \PreRelLndMax{} nodes before its publication, and
\PlatLineMin--\PlatLineMax\% of the nodes that changed line right after a platform shipped a new LND line landed on
the line it shipped (\S\ref{sec:inference}).
\item \textbf{Upgrade dynamics.} Most operators upgrade within months; the ones that lag are small or already behind. Half of the
LND nodes on an older line upgrade within \MedianLNDAll{} days of a release; the smallest quarter of nodes takes
\MedianLNDSmall{}, and a node already more than one line behind has about five times lower odds of moving in any
week. These figures set apart a managed fleet of up to \FleetPeakNodes{} live LND
nodes~\cite{kumar2026clientdiversity}, whose schedule would otherwise dominate them (\S\ref{sec:dynamics}).
\item \textbf{Platforms set the pace.} The largest upgrades are node platforms shipping a release to all their users
at once. Dating the releases of two platforms to the minute from public records, we find that the three largest LND
upgrade waves, \WaveUmbrelMinNodes--\WaveUmbrelMaxNodes{} nodes each, came about a day after the Umbrel app store
published a new LND version, and another followed BTCPay Server's security release. Since June~2025, the 48~hours
after such releases cover \PlatTimeLND\% of the time but hold \PlatShareLND\% of LND upgrades (\S\ref{sec:waves}).
\item \textbf{Disclosure, response and exposure.} Disclosure moved few operators, and mostly those already close to
current. For the \NAdvTotal{} LND and CLN advisories whose fix boundary gossip can see, we measure how many live nodes
and how much capacity stayed exposed; the fixes typically shipped months before disclosure (median \FixToDiscMedian{}
days). For \NResp{} disclosures we test whether exposed nodes left the affected lines faster afterwards. Only CLN's
July~2026 disclosures
raised that rate, by a factor of \RROpClnRecent{} among nodes one line behind; nodes further behind moved when
BTCPay's own advisory had their operators update. No LND disclosure produced a detectable increase: two months after
Lightning Labs' August~2026 batch, \PctCohortStillLLBatch\% of the nodes exposed at disclosure still ran an affected
release. When CLN asked operators who could not upgrade to restart offline, the nodes that went offline were mostly
already patched: \OffCallNewestLeft\% of the nodes on the newest line went offline, against \OffCallBehindLeft\% of
the nodes behind (\S\ref{sec:exposure}).
\end{itemize}

All data behind the figures and tables, and the code, are released. We report aggregates only: no list, ranking or
label of individual nodes by release, platform or exposure.
